\section{Skill Is Not 1 Thing: Discover Applied Capability}

\subsection{The 8 Dimensions of Real Skill}

Most founders describe skill by broad category: \textit{I know Python}, \textit{I know finance}, \textit{I am good at teaching}. These statements are too coarse to guide opportunity selection.

Skills have multiple operational dimensions:

\begin{itemize}[leftmargin=11pt,itemsep=1.5pt,topsep=1.5pt,parsep=0pt]
  \item \textbf{Execution:} Can you perform a task to the relevant quality standard?
  \item \textbf{Diagnosis:} Can you tell what is wrong when the task does not follow the textbook?
  \item \textbf{Adaptation:} Can you apply knowledge to unfamiliar cases?
  \item \textbf{Speed:} How long does reliable performance take?
  \item \textbf{Independence:} How much help, infrastructure, or supervision do you require?
  \item \textbf{Communication:} Can you make your judgment understandable to buyers and colleagues?
  \item \textbf{Repeatability:} Does the performance survive different clients and difficult conditions?
  \item \textbf{Economics:} Does the result justify a price higher than the cost of providing it?
\end{itemize}

An engineer who produces a polished prototype with an AI assistant may not be able to architect a secure, maintainable production system. A teacher who speaks eloquently may not improve student mastery. A salesperson who converts warm inbound leads at a famous brand may fail to generate demand for an unknown startup.

\subsection{Build a Task-Level Skill Map}

\textbf{Write down tasks, not disciplines.} For an AI engineer, a task-level map might include:

\begin{tcolorbox}[
  enhanced,
  colback=white,
  colframe=BayesBorder!80!white,
  boxrule=0.4pt,
  arc=2pt,
  left=3pt,right=3pt,top=3pt,bottom=3pt,
  before skip=3pt,after skip=5pt
]
\scriptsize
\setlength{\tabcolsep}{2.5pt}
\renewcommand{\arraystretch}{1.18}
\begin{tabularx}{\linewidth}{@{} >{\bfseries\color{BayesInk}}p{2.1cm} >{\color{BayesInk}}X >{\color{BayesSecondary}}p{2.3cm} @{}}
\toprule
\textbf{Task Capability} & \textbf{Unit of Assessment} & \textbf{Verified Artifact} \\
\midrule
LLM Evaluation & Design an eval suite for a retrieval workflow & Test cases, failure taxonomy, latency metrics. \\
Backend Systems & Deploy service with auth, unit tests, and metrics & Running code, tests, post-mortem records. \\
Applied Statistics & Detect selection bias in a business experiment & Diagnostic report on unseen noisy dataset. \\
Technical Pedagogy & Diagnose misconception and lift test score & Pre- and post-class blind student exams. \\
Product Scoping & Map buying journey and cut low-value features & Measured retention shift or decision memo. \\
\bottomrule
\end{tabularx}
\end{tcolorbox}

For each task, record 3 different levels:
\begin{enumerate}[label=\textbf{\arabic*.},leftmargin=11pt,itemsep=1.5pt,topsep=1.5pt,parsep=0pt]
  \item \textbf{I can explain it.} This establishes conceptual familiarity.
  \item \textbf{I can do it in a controlled environment.} This establishes applied competence.
  \item \textbf{I can deliver it for someone with real stakes.} This begins to establish commercial capability.
\end{enumerate}

\textbf{Do not grant yourself the 3rd level for having reached the 1st.}

\subsection{The Independent Work-Sample Test}

Choose 3 important capabilities from your ledger. For each, design a 2 to 6 hour trial that resembles a real buyer problem:

\begin{enumerate}[label=\textbf{\arabic*.},leftmargin=11pt,itemsep=1.5pt,topsep=1.5pt,parsep=0pt]
  \item Define an unfamiliar but realistic task.
  \item Set quality criteria before performing it.
  \item Limit time and resources to conditions a small business could afford.
  \item Produce a complete artifact, not just a narrative about how you would solve it.
  \item Ask a qualified reviewer to evaluate it without being told how confident you feel.
  \item Record the failures, rework, hours spent, and help required.
  \item Repeat once with a different task if the capability is central to the proposed venture.
\end{enumerate}

These work samples are not predictive psychological tests. Their value is practical realism: they replace vague self-belief with observable performance. They also expose hidden dependencies. If your demo requires 15 hours of unpaid customization to earn ₹5,000 of revenue, you may have a skill but not yet a viable service.

\subsection{The Skill-to-Economic-Value Bridge}

For each promising skill, complete this sentence:

\begin{tcolorbox}[
  enhanced,
  colback=BayesPale!25!white,
  colframe=BayesBorder!60!white,
  boxrule=0.4pt,
  arc=2pt,
  left=6pt,right=6pt,top=4pt,bottom=4pt,
  before skip=3pt,after skip=4pt
]
\footnotesize\color{BayesInk}
I can \textbf{[perform a specific task]} for \textbf{[a particular buyer or user]}, reducing \textbf{[cost, time, risk, or frustration]} or increasing \textbf{[revenue, quality, capability, or certainty]}, with evidence from \textbf{[an actual result]}.
\end{tcolorbox}

Example: \textit{I can diagnose poorly performing internal search systems for mid-market software teams, reducing irrelevant results and operational troubleshooting, with evidence from 2 previous deployments.}

The sentence is not a sales pitch; it is a hypothesis about economic transfer. \textbf{If you cannot identify the paying buyer and the consequential result, you may have discovered a personal talent but not yet an entrepreneurial advantage.}
